\documentclass[10pt,a4paper]{amsart}
\usepackage[margin=19mm]{geometry}
\usepackage{amsmath,amssymb,mathtools}
\usepackage[scr=rsfs,cal=euler]{mathalfa}
\usepackage{xfrac}
\usepackage[unicode,hidelinks,bookmarks=false]{hyperref}
\newcommand{\ie}{i.e.\ }
\newcommand{\cf}{cf.\ }
\newcommand{\resp}{respectively\ }
\newcommand{\Subsection}{\subsection}
\providecommand{\nobreakdash}{\nobreak\hbox{-}}
\setlength{\emergencystretch}{2em}
\multlinegap0pt
\usepackage{amssymb}
\newtheorem{theorem}{Theorem}
\title{Generalized functions\\ in a Non-Archimedean field}
\author{Vieri Benci}
\date{Source: arXiv:2609.07279v1 (CC BY 4.0)}
\begin{document}
\maketitle

\noindent\emph{Source and licence.} Generalized functions\\ in a Non-Archimedean field. Version arXiv:2609.07279v1 (CC BY 4.0), distributed by arXiv under the Creative Commons Attribution 4.0 International licence, http://creativecommons.org/licenses/by/4.0/, as recorded in the arXiv licence metadata for this exact version and shown on the abstract page https://arxiv.org/abs/2609.07279v1. Full licence notice: \texttt{licenses/arxiv-2609.07279-CC-BY-4.0.txt}.\par\smallskip

\noindent\textit{Editorial scope.} Counted unit: the article's theorem TFC (source0.tex lines 2261--2266). The article's complete original proof of it is delivered with it (source0.tex lines 2267--2282, one \texttt{proof} environment of 16 lines). No mathematics is written, repaired or re-derived: the statement and the proof are the article's own lines, in the article's own notation, and no retained line is substituted or re-flowed.  Retained uncounted as the source-local context the proof needs: lines 2109--2111.  Nothing was omitted from any retained span, so the package carries no omission record.  No retained line contains a citation, so the wrapper carries no bibliography and no bibliography entry.  No retained line contains a figure environment, an image-inclusion command or drawing code, so no image is delivered and the package declares no asset dependency.  Editorial formatting only, in addition: an AMS \texttt{amsart} preamble that restates the article's own declarations and macros used by the retained lines, namely the environments \texttt{theorem} on separate counters, together with the standard packages those lines need.  The retained environments print in this document as Theorem 1 in order of appearance, whereas the article prints the same items with the numbering of its own full document.  The complete native source of the article as distributed in the arXiv e-print for this exact version is bundled unchanged as \texttt{sources/209616/source0.tex}.
\begin{equation}
D\theta^{\circ}_{(a,b)}=\delta_{a}-\delta_{b}.\label{gru}
\end{equation}

\begin{theorem}[Fundamental Theorem of Calculus]\label{TFC} If $a,b\in\mathbb{R}$,
then for every $v\in V(\Gamma)$ 
\[
\int^{\circ}_{(a,b)}Dv\,dx=v(b)-v(a).
\]
\end{theorem}

\begin{proof}
For every $v\in V(\Gamma)$, since $\theta_{(a,b)}\in V(\mathbb{R})$,
by (\ref{gru}) we have 
\[
\int^{\circ}_{(a,b)}Dv(x)\,dx=\int^{\circ}\theta^{\circ}_{(a,b)}Dv\,dx=-\int^{\circ}D\theta^{\circ}_{(a,b)}v\,dx=-\langle\partial\theta_{(a,b)},v_{\Lambda}\rangle.
\]

Since $\partial\theta_{(a,b)}=\delta_a-\delta_b$, we get
\[
\int_{(a,b)}^{\circ} Dv(x)\,dx
=-\langle \delta_a-\delta_b, v_\Lambda\rangle
=
v(b)-v(a).
\]
$\square$
\end{proof}

\end{document}
